\section{A general bound for more than three sectors}\label{four:section}
The same comparison and witness arguments extend to any fixed number of
sectors. The resulting bound is sharp when there are three sectors, by the
main theorem, but no sharpness assertion is made for more than three.
This extension concerns the same pointwise protection of original core
maxima; it imposes no strategic reporting or incentive requirement.

Let $K\ge3$ and let $N_1,\ldots,N_K$ be disjoint sectors, each of size
$n\ge1$. For each unordered pair of distinct sectors fix a nonnegative
real pair-value matrix. A productive full team contains one player from
every sector, and its worth is the sum of all its pair values. Coalition
worth is the largest total worth of player-disjoint full teams it contains.
All smaller teams have worth zero. Assume $T=v(N)>0$ and that the entire
original core $\core$ is nonempty. Define $H_i=\max_{x\in\core}x_i$ and
\[
 s_K(M)=\min\left\{\sum_i y_i-T:
       y_i\ge H_i/K\ (i\in N),\quad y(e)\ge v(e)
       \text{ for every full team }e\right\}.
\]
As in the three-sector model, these constraints imply nonnegativity and
coverage of every original coalition.

\begin{theorem}[General sector-count bound]\label{four:main}
Every finite market in the preceding class satisfies
\begin{equation}\label{four:universal}
             s_K(M)<\frac{2(K-2)}{K(K-1)}\,T.
\end{equation}
For an integral partial matching $P$ leaving $k=n-|P|\ge1$ unmatched
players in each sector, put
\[
 J_K(M,P)=v(P)+\frac1K\sum_{i\notin\cup P}H_i-T.
\]
Then
\begin{equation}\label{four:integral}
 J_K(M,P)\le
       \frac{2(K-2)(k-1)}{K\bigl(k(K-1)-1\bigr)}\,T.
\end{equation}
If $k=0$, then $J_K(M,P)\le0$.
\end{theorem}

We give the proof because the fractional-packing step and the use of
separate witnesses are essential. Choose an optimal integral matching,
complete it to a perfect matching, and relabel its full teams by
$j=0,\ldots,n-1$. Completion preserves optimality: the sector sizes are
equal, and the added full teams have nonnegative worth. Write $d_j$ for
the worth of diagonal team $j$, and $u_{a,j}$ for the sum of its pair
values incident with its sector-$a$ player. Thus
\begin{equation}\label{four:budgets}
       \sum_j d_j=T,\qquad \sum_{a,j}u_{a,j}=2T.
\end{equation}
Every diagonal team is tight at every $x\in\core$. In particular,
\[
   \sum_{a=1}^Kx_{a,j}=d_j,\qquad
   0\le H_{a,j}\le d_j,\qquad
   \delta_{a,j}=d_j-H_{a,j}\ge0.
\]
Replacing player $(a,j)$ in its diagonal team by $(a,h)$ leaves all
pairs not incident with sector $a$ intact. Nonnegativity of the new
pairs therefore gives
\begin{equation}\label{four:range}
             x_{a,j}-x_{a,h}\le u_{a,j}
             \qquad(x\in\core).
\end{equation}
There is no division by $d_j$, so zero diagonal values are allowed.

\begin{lemma}[Aggregate comparisons for $K$ sectors]\label{four:aggregate}
Let $\lambda$ be any fractional partial matching, with
\[
 p_i=\sum_{e\ni i}\lambda_e\le1,\quad q_i=1-p_i,\quad
 W=\sum_e\lambda_ev(e),\quad S=\sum_iq_iH_i,\quad V_K=KW+S.
\]
Set
\[
 k=n-\sum_e\lambda_e,\qquad
 \beta=\sum_iq_iu_i,\qquad C_2=\sum_iq_i^2u_i,
\]
where $u_i=u_{a,j}$ for $i=(a,j)$, and put
\[
 E=\sum_{a,j}\left(\sum_{b\ne a}q_{b,j}\right)\delta_{a,j}\ge0.
\]
Then
\begin{align}
 V_K&\le (3K-4)T-(K-2)\beta-\frac{E}{K-1},
        \label{four:first}\\
 V_K&\le KT+\beta-\frac{C_2}{k}\qquad(k>0).
        \label{four:second}
\end{align}
\end{lemma}
\begin{proof}
For a full team $e=((1,j_1),\ldots,(K,j_K))$ and distinct sectors
$b,c$, subtract diagonal tightness at $j_b$ from coverage of $e$.
For every $x\in\core$, apply~\eqref{four:range} to the other $K-2$
sectors to obtain
\[
 x_{c,j_c}-x_{c,j_b}
 \ge v(e)-d_{j_b}-\sum_{a\notin\{b,c\}}u_{a,j_a}.
\]
Choose an original core point attaining the subtracted maximum
$H_{c,j_b}$, while bounding its other coordinate by $H_{c,j_c}$.
This proves the same inequality for the difference of those two maxima.
Each ordered pair $(b,c)$ may use a different attaining core point.
Sum over all $K(K-1)$ ordered pairs, writing
$H_j^\Sigma=\sum_aH_{a,j}$:
\begin{equation}\label{four:team-comparison}
\begin{split}
 K(K-1)v(e)\le{}&(K-1)\sum_a d_{j_a}
       +(K-1)(K-2)\sum_a u_{a,j_a}\\
       &+K\sum_aH_{a,j_a}-\sum_aH_{j_a}^\Sigma.
\end{split}
\end{equation}
Here every incident budget occurs $(K-1)(K-2)$ times, and the ideal
terms are $K\sum_aH_{a,j_a}-\sum_aH_{j_a}^{\Sigma}$.
The count preserves multiplicities when some team labels coincide.

Multiply~\eqref{four:team-comparison} by $\lambda_e$ and sum over
teams. Add $(K-1)S$. If $n_j=\sum_bp_{b,j}$, the resulting coefficient
of $H_{a,j}$ is
\[
       Kp_{a,j}+(K-1)q_{a,j}-n_j
                   =\sum_{b\ne a}q_{b,j}.
\]
Replacing $H_{a,j}$ by $d_j-\delta_{a,j}$ gives an ideal contribution
$(K-1)\sum_{a,j}q_{a,j}d_j-E$. The other diagonal contribution is
$(K-1)\sum_{a,j}p_{a,j}d_j$. These sum to $K(K-1)T-E$.
The budget contribution is
$(K-1)(K-2)(2T-\beta)$ by~\eqref{four:budgets}.
The left-hand side is $(K-1)V_K$, proving~\eqref{four:first}.

Every sector has residual mass $k$: each full team uses exactly one
player per sector. For every original core point $x$, fractional team
coverage gives $\sum_i p_ix_i\ge W$, and hence
\begin{equation}\label{four:residual}
                   \sum_iq_ix_i\le T-W.
\end{equation}
In particular $W\le T$. Fix a player $i$ and choose a core point
attaining $H_i$. Its other same-sector coordinates are at least
$H_i-u_i$ by~\eqref{four:range}, while all other-sector coordinates
are nonnegative. Therefore
\[
        T-W\ge q_iH_i+(k-q_i)(H_i-u_i),
        \qquad
        kH_i\le T-W+(k-q_i)u_i.
\]
The factor $k-q_i$ is the sum of the other residual weights in $i$'s
sector and is nonnegative. Multiplying by $q_i$, summing the resulting
scalar inequalities, and using $\sum_iq_i=Kk$ yields
\[
                    kS\le Kk(T-W)+k\beta-C_2.
\]
For $k>0$, division by $k$ proves~\eqref{four:second}. Separate
maximizing witnesses suffice because only the fixed numbers $H_i$
appear in these scalar inequalities.
\end{proof}

\begin{proof}[Proof of Theorem~\ref{four:main}]
The shifted covering-program dual is
\begin{equation}\label{four:dual}
 s_K(M)=\max_{\lambda\in\pack_K}\left(W+\frac SK-T\right)
       =\max_{\lambda\in\pack_K}\left(\frac{V_K}{K}-T\right),
\end{equation}
where $\pack_K$ is the fractional partial-matching polytope just defined.
Indeed, substitute $y_i=H_i/K+z_i$, $z_i\ge0$, in the covering program;
its dual has the load constraints $p_i\le1$. Adding the constant
$\sum_iH_i/K-T$ gives~\eqref{four:dual}. The primal is feasible and
bounded below, while $\pack_K$ is compact, so both optima are attained.

For $k>0$, take weights $1/(K-1)$ and $(K-2)/(K-1)$ on
\eqref{four:first} and~\eqref{four:second}, respectively. The
$\beta$ terms cancel, giving
\begin{equation}\label{four:fractional-combination}
 V_K\le\left(K+\frac{2(K-2)}{K-1}\right)T
         -\frac{E}{(K-1)^2}
         -\frac{(K-2)C_2}{(K-1)k}.
\end{equation}
If $C_2>0$, this is strictly smaller than the displayed leading term.
If $C_2=0$, nonnegativity of each $q_i^2u_i$ implies $\beta=0$, so
\eqref{four:second} gives $V_K\le KT$, again strictly smaller since
$K\ge3$ and $T>0$. When $k=0$, all residual weights vanish and
\eqref{four:residual} gives $W\le T$, so $V_K\le KT$ as well.
Apply this strict bound at a maximizing point of the compact dual
polytope to prove~\eqref{four:universal}. A pointwise strict inequality
without attained maximization would not by itself justify this step.

For an integral partial matching with $k\ge1$, residual weights belong
to $\{0,1\}$ and $C_2=\beta$. Drop $E\ge0$ from
\eqref{four:first} and take weights
\[
       \frac{k-1}{k(K-1)-1},\qquad
       \frac{k(K-2)}{k(K-1)-1}
\]
on~\eqref{four:first} and~\eqref{four:second}. They are nonnegative
and sum to one. The $\beta$ terms again cancel, giving
\[
       V_K\le\left(K+
           \frac{2(K-2)(k-1)}{k(K-1)-1}\right)T.
\]
Since $J_K=V_K/K-T$, this proves~\eqref{four:integral}, including
$k=1$. For $k=0$, the matching is perfect and $J_K=v(P)-T\le0$.
\end{proof}

\begin{remark}[Four sectors and the scope of sharpness]\label{four:scope}
For $K=4$, the theorem gives $s_4(M)<T/3$ and
$J_4(M,P)\le(k-1)T/(3k-1)$. It therefore gives a supremum upper bound
$\sup_M s_4(M)/T\le1/3$; it does not prove that this supremum equals
$1/3$, nor a smaller fixed gap below $1/3$ uniform over market sizes.
The analogous distinction applies for every $K>3$. The integral bound
is an upper bound on one integral dual obstruction, not on the full
fractional optimum. When $K=3$, both formulas recover the sharp bounds
already proved in the main text. Sparse allowed sector graphs are
included by setting the unused pair matrices to zero.
\end{remark}

\begin{remark}[What the exchange budgets provide]\label{four:exchange}
More generally, the same argument works for full-team games with a tight
optimal block partition and nonnegative numbers $u_i$ satisfying the
directed same-sector range bound~\eqref{four:range}. If
$B=\sum_i u_i$, the first bound becomes
\[
       V_K\le KT+(K-2)(B-\beta)-E/(K-1),
\]
and the same weighted combination gives
$s_K(M)\le (K-2)B/[K(K-1)]$.
Pair nonnegativity supplies these range bounds through an actual
replacement-team calculation, and pair additivity supplies $B=2T$.
These hypotheses cannot be replaced merely by assumed versions of the
two aggregate inequalities. This is a direct extension of the proof,
not a separately established sharpness or historical-priority claim.
\end{remark}

\section{Protected ideals in short flow and transportation games}\label{flow:section}
This appendix concerns a different cooperative technology from the
pair-additive assignment model. It tests which protection statements extend
to arc-player flow games. The results are pointwise statements about each
original entire core. No incentive guarantee across capacity reports is
asserted. In particular, the counterexample below refutes simultaneous
endpoint-ideal attainment while having zero protected subsidy itself.

\subsection{The model and the sharp unit-capacity bound}
Let $G$ be a finite directed acyclic network with distinct source $s$ and
sink $t$. Each arc is owned by a distinct player; parallel arcs are
permitted. A coalition's worth is its maximum $s$--$t$ flow using only its
arcs. Write $T=v(N)>0$, let $\core$ be the entire original core, and set
$H_e=\max_{x\in\core}x_e$. For this appendix define
\[
 s_{1/3}(G)=\min\{y(N)-T:y\ge H/3,\ y(A)\ge v(A)
                                   \text{ for every }A\subseteq N\}.
\]
For any minimum cut, paying each of its arcs its capacity and all other
arcs zero gives a core allocation: its total is $T$, and each coalition's
flow is at most the capacity of its available arcs in that cut. This
also proves nonemptiness of the original core. The core is compact, so
all the maxima defining $H$ are attained.

\begin{lemma}[Deletion and unit ideals]\label{flow:ideals}
For every grand maximum flow $f$ and arc $e$,
\[
                    H_e\le T-v(N\setminus\{e\})\le f_e.
\]
If all capacities are one, $H_e\in\{0,1\}$, and $H_e=1$ precisely
when $e$ belongs to some minimum cut.
\end{lemma}
\begin{proof}
The first inequality follows by combining core coverage of
$N\setminus\{e\}$ with efficiency. Decompose $f$ into paths and remove
all paths using $e$; the lost value is $f_e$, giving the second inequality.
With unit capacities, integral maximum flows show that the deletion
marginal is zero or one. If it is one, a minimum cut after deleting $e$
has capacity $T-1$. The restored arc must cross that cut and makes it a
minimum cut of the original network. Its cut allocation attains payoff
one at $e$. Conversely, membership in a minimum cut gives that same
attainment. If the deletion marginal is zero, so is $H_e$.
\end{proof}
The unit-ideal characterization, cut allocations, and minimum-cut lattice
are classical flow-game ingredients~\cite{flow:Fang2010,flow:KalaiZemel1980}.
They are included to separate them from the protected-subsidy conclusion.

\begin{theorem}[Short unit-capacity flows]\label{flow:unit}
Suppose every directed $s$--$t$ path has at most three arcs, and all arc
capacities are one. Then
\begin{equation}\label{flow:unit-bound}
                     s_{1/3}(G)\le\frac{T-1}{3}.
\end{equation}
For every positive integer $T$, equality occurs in a network whose
$s$--$t$ paths all have three arcs. Consequently the supremum of
$s_{1/3}(G)/T$ in this class is $1/3$.
\end{theorem}
\begin{proof}[Upper bound]
Directed cut capacity is submodular in the source-side vertex set.
The intersection and union of two minimum-cut source sides are therefore
also minimum-cut source sides: both crossed cuts have capacity at least
$T$, while submodularity bounds their sum by $2T$. Since the network is
finite, the intersection $U_-$ and union $U_+$ of all minimum-cut source
sides define minimum cuts $C_-$ and $C_+$.

Every source arc $(s,u)$ with ideal one crosses some minimum cut, so $u$
is missing from at least one source side and hence from $U_-$. Thus
$C_-$ simultaneously contains every positive-ideal source arc. Similarly,
$C_+$ contains every positive-ideal sink arc $(u,t)$. Direct $s$--$t$
arcs cause no difficulty and can be in both cuts.

Let $B$ be the positive-ideal arcs outside $C_-\cup C_+$. Such arcs are
strictly internal. Fix an integral grand maximum flow decomposed into
$T$ arc-disjoint paths. Lemma~\ref{flow:ideals} implies that every arc
in $B$ is used by this flow. Each path has at most one strictly internal
arc, so $r:=|B|\le T$. If $r=0$, the average of the two cut allocations
already dominates $H/3$ and has zero subsidy. If $r>0$, choose $e\in B$
and a minimum cut $C_e$ containing it. Average the allocations of
$C_-,C_+,C_e$, and raise any remaining coordinates to their floors $H/3$.
Only the at most $r-1$ positive-ideal arcs outside these three cuts can
need a raise, of at most $1/3$ each. Raising coordinates preserves every
coalition inequality. This proves~\eqref{flow:unit-bound}.
\end{proof}

\begin{proposition}[A sharp chain family]\label{flow:chain}
For $n\ge1$, use vertices $s,L_1,\ldots,L_n,R_1,\ldots,R_n,t$ and
unit-capacity arcs
\[
 a_i=sL_i,\quad b_i=L_iR_i,\quad c_i=R_it\quad(1\le i\le n),
 \qquad d_i=L_iR_{i+1}\quad(1\le i<n).
\]
Then $T=n$, the ideals of the $3n$ arcs $a_i,b_i,c_i$ are one, all
$d_i$ ideals are zero, and $s_{1/3}(G)=(n-1)/3$.
\end{proposition}
\begin{proof}
The diagonal paths attain flow $n$, matching the source cut. Their
maximum flow avoids the $d_i$, so Lemma~\ref{flow:ideals} gives
$H_{d_i}=0$. The source and sink cuts attain all $a_i$ and $c_i$ ideals.
For $b_i$, the vertex set
\[
           \{s,L_i,\ldots,L_n,R_{i+1},\ldots,R_n\}
\]
has boundary $\{a_j:j<i\}\cup\{b_i\}\cup\{c_j:j>i\}$, a minimum
cut of size $n$. Hence all $b_i$ ideals are also one.

The $n-1$ cross paths $(a_i,d_i,c_{i+1})$ are arc-disjoint. Their union
has worth $n-1$ and leaves $n+2$ positive-ideal arcs outside it. Its
coalition inequality and those exterior floors force protected total at
least $n-1+(n+2)/3$, giving subsidy at least $(n-1)/3$.
For the reverse bound pay $1/3$ to every $a_i,b_i,c_i$, an additional
$1/3$ to $a_i$ for $i<n$, and zero to $d_i$. Every diagonal or cross
path receives at least one. Every coalition's unit-capacity integral
maximum flow decomposes into arc-disjoint such paths, so summing path
inequalities covers every coalition. The total is $n+(n-1)/3$.
\end{proof}
This proves the equality and supremum statements in
Theorem~\ref{flow:unit}. The direct path-cover argument just used depends
on unit capacities; it is not a proof for arbitrary capacities.

\subsection{Prerequisite transfers and a one-sided quota extension}
\begin{lemma}[Prerequisite transfer]\label{flow:transfer}
Let $v$ be a nonnegative monotone game with a nonempty nonnegative core.
Suppose distinct players $a,e$ satisfy
$v(A)=v(A\setminus\{e\})$ whenever $a\notin A$. Transferring all of
$e$'s payoff to $a$ preserves core membership.
\end{lemma}
\begin{proof}
Efficiency and nonnegativity are preserved. Coalitions containing both
or neither player have unchanged payoff, and coalitions containing $a$
only gain. For a coalition containing $e$ but not $a$, the new payoff is
$x(A\setminus\{e\})\ge v(A\setminus\{e\})=v(A)$.
\end{proof}

In a pure three-layer network $s\to L\to R\to t$, let $M$ be the
middle-arc players and $D$ the endpoint-arc players. Define the endpoint
game $w(Q)=v(Q\cup M)$ for $Q\subseteq D$. Its entire core describes
exactly the zero-middle face:
\begin{equation}\label{flow:face}
 \core\cap\{x_M=0\}=\{(z,0_M):z\in\core(w)\}.
\end{equation}
To see this, restrict by testing coalitions $Q\cup M$; conversely, for
any original coalition $A$ use
$v(A)\le w(A\cap D)$. Grand values agree. If every $R_j$ has a unique
sink arc, each middle arc ending there requires that sink arc on any
productive path. Lemma~\ref{flow:transfer} moves all middle payments
sinkward, leaves every source payoff unchanged, and weakly increases
every sink payoff. Thus the source projections of the full core and
this face coincide, and every endpoint coordinate ideal is preserved.
Unique source arcs give the symmetric conclusion by sourceward transfer.
Neither statement changes any original owner or capacity. If both
endpoint ports have parallel owners, neither map is automatically valid.

\begin{theorem}[One-sided quotas]\label{flow:quota}
Fix a rational $q>0$. Consider a pure three-layer network with one or more source
arcs into each $L$ vertex, source capacities $q r_i$ for integers
$r_i\ge0$, exactly one sink arc of capacity $q$ out of each $R$ vertex,
and each rational middle capacity either zero or at least $q$. Parallel source
and middle arcs are allowed; each arc remains a distinct player.
All source ideals are jointly attained in the original core, and all
sink ideals are jointly attained at another core point. If $m$ middle
arcs have positive ideal, then
\begin{equation}\label{flow:quota-bound}
 s_{1/3}(G)\le\frac q3\max\{0,m-1\}\le\frac{T-q}{3}.
\end{equation}
This last bound is attained for every positive integer $T/q$.
\end{theorem}
\begin{proof}
Scale to $q=1$. Because the only arc out of each $R$ vertex has capacity
one, capping every positive middle capacity at one preserves every
coalition value. Zero-capacity arcs have zero core payoff. The normalized
network is integral, so $T$ is a positive integer. The sinkward transfer
and~\eqref{flow:face} preserve all endpoint ideals.

With all middle arcs available, treat each source arc as a firm with
quota $r_i$ and each sink arc as a unit worker. Their assignment value
$a_{ij}$ is one if a positive middle arc links their endpoints, and zero
otherwise. Integral flows give matchings of this many-to-one market.
Conversely, each allowed firm--worker match can be routed through an
available middle arc; a worker is used at most once, so no middle arc
of capacity one restricts the matching. Distinct source arcs sharing a
head are distinct firms with the same allowable workers, not clones
introduced by the proof.

For completeness, we give the standard salary-lattice argument behind
side-optimality~\cite{flow:AtayNunezSolymosi2026}. It works for any
nonnegative assignment matrix $a$. Remove zero-quota firms, and balance
total firm quota and worker count by adding zero-value dummy workers
or a zero-value dummy firm. These are null players; their core payoffs
are zero and projecting them away preserves the original core. Complete
an optimal matching $\mu$ to fill every quota and assign every worker.
The added matches have value zero. Let $\mu(i)$ be the workers matched
to firm $i$, and $i(j)$ worker $j$'s matched firm.

A payoff vector $(x,y)$ is in the balanced core precisely when
\begin{align}
 0&\le y_j\le a_{i(j),j},\label{flow:salary-box}\\
 y_k-y_j&\ge a_{ik}-a_{ij}
         &&(j\in\mu(i),\ k\notin\mu(i)),\label{flow:salary-diff}\\
 x_i&=\sum_{j\in\mu(i)}(a_{ij}-y_j).\label{flow:salary-firm}
\end{align}
Indeed, the disjoint coalitions $\{i\}\cup\mu(i)$ attain the grand
value, so their core inequalities are all tight, giving
\eqref{flow:salary-firm}. Deleting one matched worker gives the upper
bound in~\eqref{flow:salary-box}; replacing that worker by an unmatched
worker gives~\eqref{flow:salary-diff}. These use feasible matching
lower bounds, without asserting that a replacement is itself optimal.
Conversely, for any firm $i$ and worker set $U$ of size at most $r_i$,
pair each $k\in U\setminus\mu(i)$ with a distinct
$j\in\mu(i)\setminus U$. Equation~\eqref{flow:salary-diff} gives
$a_{ik}-y_k\le a_{ij}-y_j$, and all remaining matched surpluses are
nonnegative by~\eqref{flow:salary-box}. Hence
\[
                   x_i\ge\sum_{k\in U}(a_{ik}-y_k).
\]
Every feasible coalition matching partitions into disjoint one-firm
teams; summing these inequalities, together with nonnegative unmatched
payoffs, proves all core constraints. Efficiency follows from
\eqref{flow:salary-firm}.

The nonempty compact salary set defined by
\eqref{flow:salary-box}--\eqref{flow:salary-diff} is closed under both
coordinatewise minimum and maximum. For example, two inequalities
$y_k\ge y_j+d$ and $z_k\ge z_j+d$ imply the corresponding minimum
and maximum inequalities. Taking the minimum of finitely many
coordinate-minimizing salary vectors therefore minimizes every salary
simultaneously; by~\eqref{flow:salary-firm} it maximizes every firm's
payoff simultaneously. Coordinate maxima give a point maximizing every
worker payoff. Project away null players and lift through
\eqref{flow:face} to obtain the two asserted original-core endpoint
allocations.

For a normalized middle arc, deletion lowers the integral maximum flow
by zero or one. If it lowers it by one, a minimum cut after deletion,
with the arc restored, is an original minimum cut attaining payoff one
there. Thus each middle ideal is zero or one. Every positive-ideal
middle arc is used by every integral grand maximum flow. Each unit of
flow uses one middle arc, so $m\le T$.

If $m=0$, average the two side-optimal endpoint allocations; this already
dominates $H/3$. If $m>0$, average these two core points and a core point
attaining one positive middle ideal, then raise the other middle
coordinates to their floors. At most $m-1$ coordinates require a raise
of at most $1/3$ each. Rescale by $q$ to obtain~\eqref{flow:quota-bound}.
The chain in Proposition~\ref{flow:chain}, scaled by $q$, lies in this
class and proves sharpness.
\end{proof}
The proof allows parallel source owners because the unique sink supplies
the prerequisite. Arbitrary capacities on both endpoint sides are not
covered: then an endpoint owner may serve several opposite owners, and
the disjoint one-firm-team salary argument need not apply. Splitting
that owner into independent players would change the entire core.

\subsection{An exact transportation embedding}
\begin{proposition}[Embedding without cloning endpoint owners]\label{flow:embedding}
Let disjoint endpoint sets $I,J$ have capacities $r_i,c_j\ge0$ and
allowed pairs $E\subseteq I\times J$. A coalition's worth $w$ is the
maximum transported amount on its allowed pairs, subject to these row
and column capacities, with no separate pair-capacity restriction.
Assume $T=w(I\cup J)>0$. There is a three-layer arc-player flow game
whose entire core is exactly
\[
                  \{(x,0_D):x\in\core(w)\},
\]
where endpoint players retain their identities and $D$ consists of new
middle-arc owners.
\end{proposition}
\begin{proof}
Set $M=\min\{\sum_i r_i,\sum_j c_j\}$. Use source arc $sL_i$ of
capacity $r_i$, sink arc $R_jt$ of capacity $c_j$, and for every allowed
pair two parallel middle arcs $L_iR_j$, each of capacity $M$ and with
its own new owner. With every middle arc available, their capacities
cannot restrict any transportation plan, so endpoint coalitions with
all middle owners have worth $w$.
Deleting any one middle arc leaves its twin of capacity $M$, and hence
does not change the grand value. Every such arc therefore has payoff
zero at every original core point by the deletion bound. Testing endpoint
coalitions together with all middle owners shows that the endpoint
restriction belongs to $\core(w)$. Conversely, extend any point of
$\core(w)$ by zero. The value of any flow coalition is at most the
transportation value of its endpoint restriction, so every original
coalition is covered. Efficiency gives the claimed equality of cores.
\end{proof}
These new middle owners need not be null players of the characteristic
function: they can help coalitions lacking their twins. Only their
zero payoff in every grand-core allocation is used. All returns per
transported unit are one; weighted transportation and a one-unit
constraint per pair are different models.

\subsection{Endpoint ideals need not be jointly attainable}
\begin{theorem}[Twelve-owner endpoint counterexample]\label{flow:endpoint}
In the two-hop network $s\to u\to t$, give the five source arcs
$P,A,G,M,B$ capacities
\[
                         (2,12,46,62,75)
\]
and the seven sink arcs $Q,R,D,E,F,L,N$ capacities
\[
                         (6,10,16,22,24,56,62).
\]
Every arc has a separate owner. Then $T=196$, and the entire-core ideals
are
\[
 H_{\rm source}=(1,6,23,31,38),\qquad
 H_{\rm sink}=(6,10,16,22,24,56,62).
\]
The largest total source payoff is $98$, although the source ideals sum
to $99$. Already $\max_{x\in\core}(x_A+x_B)=43< H_A+H_B=44$.
Nevertheless $s_{1/3}(G)=0$.
\end{theorem}
\begin{proof}
The source capacity sum is $197$ and the sink capacity sum is $196$.
For source set $S$ and sink set $R$, the coalition value is
$\min\{c(S),c(R)\}$. Introduce baseline weights
\[
 a=(1,6,23,31,37),\qquad b=(3,5,8,11,12,28,31),
                  \qquad a(\mathrm{sources})=b(\mathrm{sinks})=98.
\]
Source capacities are $2a$, except that $B$ has one extra unit; sink
capacities are $2b$. We will prove that
\begin{equation}\label{flow:UV}
         U=(a;b),\qquad V=(1,5,23,31,38;b)
\end{equation}
are core allocations, and then prove the matching real-valued upper
bounds on all individual and joint source payoffs.

Let $C_0$ denote the ten players other than $A,B$. Their baseline
weights are $(1,3,5,8,11,12,23,28,31,31)$, of total $153$.
They have no subset of weight $61$, and therefore none of weight $92$.
Here is a direct check. The smaller weights
$D_0=(1,3,5,8,11,12)$, of total $40$, have no subset of weight
$2,7,10$: for $2$ only $1$ fits; for $7$, the weights $1,3,5$ do
not suffice exactly; for $10$, using $8$ leaves $2$, while the smaller
weights total $9$. A subset of weight $61$ uses at most one $31$.
With a $31$, the remaining $30$ either uses $28$ and leaves $2$, uses
$23$ and leaves $7$, or uses only $D_0$ whose complement would have
weight $10$. Without a $31$, using both $28,23$ leaves $10$; using
only $28$ leaves a $D_0$ subset of weight $33$ with complement $7$;
using only $23$ leaves weight $38$ with complement $2$; using neither
cannot reach $61$. These exhaust the cases. Consequently every subset
of all twelve baseline weights of sum $98$ contains $A,B$ either both
or neither: including only $A$ would require a $C_0$ subset of weight
$92$, and including only $B$ would require one of weight $61$.

For arbitrary $S,R$, put $d=a(S)-b(R)\in\mathbb Z$ and
$\delta=\mathbf1_{B\in S}$. If $d\le-1$, coalition worth is
$2a(S)+\delta$, and $U$'s slack is $-d-\delta\ge0$.
If $d\ge0$, worth is $2b(R)$ and its slack is $d\ge0$.
For $V$, these slacks become $-d-\mathbf1_{A\in S}$ when $d\le-1$
and $d+\mathbf1_{B\in S}-\mathbf1_{A\in S}$ when $d\ge0$.
They are nonnegative unless possibly $d=0$. In that case the subset
$S\cup(\mathrm{sinks}\setminus R)$ has baseline weight $98$,
so $A,B$ have equal membership and the slack is zero.
Thus both nonnegative total-$196$ vectors in~\eqref{flow:UV} cover
every coalition.

Now let $x$ be any original core point. For sources define loss
$z_i=c_i-x_i$, and for sinks put $t_j=x_j$. Efficiency yields
\begin{equation}\label{flow:loss-total}
                       \sum_i z_i-\sum_j t_j=1.
\end{equation}
The complement constraint for source $P$ gives $z_P\ge1$;
all source losses are nonnegative by their complement constraints.
Let $\theta$ equal $z$ on sources and $t$ on sinks, and write
\[
 p=\theta_P,\ q=\theta_Q,\ r=\theta_R,\ d_0=\theta_D,
 \ e=\theta_E,\ f=\theta_F,\ g=\theta_G,\ l=\theta_L,
 \ m=\theta_M,\ n=\theta_N,\quad
 \alpha=\theta_A,\ \beta=\theta_B.
\]
These are proof variables only; in particular $p\ge1$.

Whenever $a(S)=b(R)$, both $(S,R)$ and its complement have worth
equal to their sink capacity, and their worths sum to $196$.
Efficiency makes both core inequalities tight, so
$z(S)-t(R)=\mathbf1_{B\in S}$. Equivalently, if a subset $W$ of all
players has baseline weight $98$, then
$\theta(W)=\sum_jt_j+\mathbf1_{B\in W}$.
For every $C_0$ subset $K$ of baseline weight $55$, apply this to
$W=\{A,B\}\cup K$ and use~\eqref{flow:loss-total} to obtain
\[
                     \theta(K)=k_0:=p+g+m.
\]
The eight such subsets
\[
 PGM,\ PGN,\ PEFM,\ PQDFM,\ RDEM,\ PQGL,\ PQREFG,\ PDEFG
\]
therefore have equal $\theta$-sum. In succession, these equalities give
\[
 m=n,\quad g=e+f,\quad e=q+d_0,\quad p+f=r+d_0,
 \quad m=q+l,\quad m=q+r+g,\quad m=d_0+g.
\]
Solving them yields
\begin{align}
 d_0&=q+r,&e&=2q+r,&f&=q+2r-p,&g&=3q+3r-p,\notag\\
 l&=3q+4r-p,&m=n&=4q+4r-p,&k_0&=7q+7r-p.
       \label{flow:identities}
\end{align}
Thus $\theta(C_0)=-4p+19q+20r$. Combining this with efficiency gives
\begin{equation}\label{flow:loss-formulas}
 \alpha+\beta=1-2p+5q+6r,\qquad
                 \sum_i z_i=1-3p+12q+13r.
\end{equation}

If $W$ has baseline weight $97$, take $S=W\cap\mathrm{sources}$
and $R=\mathrm{sinks}\setminus W$. Then $a(S)-b(R)=-1$, so its
actual source capacity is strictly smaller than its sink capacity.
Core coverage gives $z(S)\le t(R)$, or
$\theta(W)\le\sum_jt_j$. Hence for every $C_0$ subset $K$ of
weight $54$,
\[
                         \theta(K)\le k_0-1.
\]
Use $K_1=\{M,E,D,Q,P\}$ and $K_2=\{L,F,D,R,P\}$.
Their $\theta$-sums from~\eqref{flow:identities} are $8q+6r$ and
$5q+8r-p$, respectively, so
\[
              r-q\ge p+1,\qquad r\le2q-1.
\]
Consequently
\begin{equation}\label{flow:qr}
                     q\ge p+2,\qquad r\ge2p+3,\qquad p\ge1.
\end{equation}
For the weight-$97$ subset $\{B,M,L,P\}$, whose $C_0$ portion has
$\theta$-sum $k_0+r$, the same strict-coalition inequality gives
$\alpha\ge r+1\ge6$. For the weight-$97$ subset
$\{A\}\cup(C_0\setminus\{M,N\})$, it gives
\[
                    \beta\ge -p+4q+5r+1\ge37.
\]
Also~\eqref{flow:identities} and~\eqref{flow:qr} imply
$g\ge8p+15\ge23$ and $m\ge11p+20\ge31$.
Thus every source payoff is coordinatewise bounded above by
$(1,6,23,31,38)$. Vector $U$ attains the first four bounds and $V$
attains the fifth. For each sink, deletion gives $H_j\le c_j$;
the all-sink cut allocation attains all those bounds at once.

Finally~\eqref{flow:loss-formulas}--\eqref{flow:qr} give
\[
             \sum_i z_i\ge35p+64\ge99,
              \qquad \alpha+\beta\ge15p+29\ge44.
\]
Source total is therefore at most $197-99=98$, and $x_A+x_B$ at
most $12+75-44=43$. Both $U,V$ attain those support bounds. This
proves the claimed incompatibility. Vector $U$ already dominates
$H/3$ on every source and sink and is efficient, so the protected
subsidy is zero.
\end{proof}

The baseline pattern is drawn from Isbell's integer majority-weight
representations~\cite[p.~27]{flow:Isbell1959}; the weights are not
claimed as a new construction ingredient. Integer nonuniqueness alone
would not prove Theorem~\ref{flow:endpoint}, which uses additional
all-coalition and real-valued core arguments. With all allowed-pair
returns one, this is also unrestricted unweighted transportation.
It matches the historical side-optimal-core question stated in
N\'u\~nez's 2013 lecture slides~\cite[pp.~37--40]{flow:Nunez2013}.
Neither its novelty relative to later or unread literature nor the
current openness of that historical question is established here.
The weighted multiple-partners model with capacity one per pair is
different and is not settled by this example.

\subsection{Why mixed tight-set uncrossing is insufficient}
For endpoint sources $I$, sinks $J$, and fixed public middle capacities,
let $w(Q,R)$ be the flow value when endpoint sets $Q\subseteq I$ and
$R\subseteq J$ are active. The established network-rank inequality is
\begin{equation}\label{flow:mixed}
 w(Q,R)+w(Q',R')\ge
 w(Q\cap Q',R\cup R')+w(Q\cup Q',R\cap R').
\end{equation}
This is the mixed submodularity of network linking
ranks~\cite{flow:Goemans2011}. A direct cut proof also fixes the direction
needed here. Use the full fixed vertex set, minimum-cut source sides
$U,V$ for the original two endpoint profiles, and source sides
$U\cap V,U\cup V$ for the two crossed profiles in the order displayed.
The public-middle cut contributions satisfy submodularity. For a source
arc, write $a,a'$ for its activity indicators and $u,v$ for the
indicators that its head is outside $U,V$. Then
\[
 au+a'v\ge (a\wedge a')\max(u,v)+(a\vee a')\min(u,v).
\]
For a sink arc use activity indicators $b,b'$ and tail-membership
indicators $u,v$; its crossed contribution is
$(b\vee b')\min(u,v)+(b\wedge b')\max(u,v)$, also at most
$bu+b'v$. Multiply by nonnegative capacities, add the middle-arc
inequality, and bound each crossed cut below by its maximum-flow value.
This proves~\eqref{flow:mixed}.

Additive payoffs have equal total payment on an original pair and its
mixed-crossed pair. The inequality~\eqref{flow:mixed} has the wrong
direction to force the crossed coalitions tight when the originals are
tight. The failure occurs even in a two-hop game: two source owners have
capacities $3,3$, three sink owners capacities $2,2,1$, and the core
allocation is $(0,0;2,2,1)$. The coalitions
\[
 (\{a_1\},\{b_1,b_3\}),\qquad
 (\{a_1,a_2\},\{b_2,b_3\})
\]
each have worth and payment $3$. Their crossed pair has worths $3,1$
but payments $5,1$. Thus total worth falls from $6$ to $4$ while
payment stays $6$, although the original coalitions overlap and neither
contains the other. This disproves mixed-crossed tight-set closure, not
every conceivable use of mixed submodularity or side-optimality.

\paragraph{Remaining scope.}
The general rational-capacity conjecture
$s_{1/3}(G)\le T/3$ for all arc-player networks of path length at most
three remains unresolved by these results. The endpoint example has
subsidy zero, and the unit and quota theorems make restricted assumptions.
For example, the serial path of capacities $(2,2,1)$ has the unit-simplex
entire core and $H=(1,1,1)$, although only its last arc lies in a minimum
cut. Clipping that cut allocation at $H/3$ costs $2/3$, whereas the
protected subsidy is zero at $(1/3,1/3,1/3)$. This illustrates why
entire-core maxima cannot be replaced by cut or optimal-dual maxima at
unequal capacities. No owner cloning, simple-graph transfer, or all-profile
truthfulness assertion is implicit in this appendix.

\section{Additional structure of frozen-ideal protection}\label{anc:section}
The results below retain structural information from the sharp-subsidy
program that is not implied by the numerical bound alone. The compact-core
representation and the finite cyclic example concern sparse supports;
the incompatibility and perturbation results are all-size statements.
They remain parts of the same protected-ideal question.

\subsection{A compact entire-core representation on sparse supports}
Suppose the three pair matrices have support contained in the pairs of
an optimal diagonal matching $\mu$ on $m$ labels and an ordinary
player-disjoint partial matching $P$ of $r$ full teams. A pair is called
supported if it appears in $\mu$ or $P$, even when its value is zero.
Teams of $P$ may repeat diagonal labels or share pairs with $\mu$.
Let $T$ be the diagonal matching's worth and assume the entire core is
nonempty.

\begin{proposition}[Sector-minimum representation]\label{anc:compression}
A vector $x$ is in the entire original core if and only if there exist
three variables $\ell_A,\ell_B,\ell_C$ such that
\begin{enumerate}
\item $x\ge0$, $\ell_s\ge0$, and $\ell_s\le x_i$ for every sector-$s$
player $i$;
\item for each supported pair $\{i,j\}$ of value $a_{ij}$, whose missing
sector is $u$, $x_i+x_j+\ell_u\ge a_{ij}$;
\item the full-team inequalities hold for every team of $\mu\cup P$;
\item for each supported off-diagonal pair joining labels $a\ne b$, its
two full-team inequalities obtained by completing it at label $a$ or
label $b$ hold;
\item $x(N)=T$.
\end{enumerate}
The system uses $3m+3$ variables, at most $7m+10r$ inequalities apart
from variable nonnegativity, and one total equality. No equality forcing
$\ell_s$ to be a coordinate minimum is needed.
\end{proposition}
\begin{proof}
Given a core point, take the actual sector minima for $\ell$.
Completing a supported pair with a player attaining the missing-sector
minimum gives the pair inequality, because all other pair values are
nonnegative. The listed full-team constraints and efficiency also hold.

Conversely, consider an arbitrary full team. If it has two supported
off-diagonal pairs, their common player belongs to one team of the
actual partial matching $P$, so both pairs belong to that same team.
Their union is therefore a team of $P$ already listed. If it has one
supported off-diagonal pair, a third label equal to one endpoint label
gives a listed local completion. Otherwise there is no diagonal pair,
the remaining off-diagonal pairs are unsupported, and its worth equals
that single pair value. The pair inequality and $\ell_u\le x_k$
then give its full-team inequality. With no supported off-diagonal
pair, the team is either diagonal, or has at most one supported diagonal
pair, or has worth zero. The listed diagonal inequality, pair-plus-minimum
inequality, or nonnegativity covers these cases. Summing over disjoint
full teams proves every coalition constraint.

There are $3m$ lower-envelope inequalities, at most $3m+3r$ pair
inequalities, $m+r$ diagonal or $P$ team inequalities, and at most
$6r$ completion inequalities. Their total is $7m+10r$; duplicates can
only reduce it.
\end{proof}
For example, a supported off-diagonal $AB$ pair with labels $a\ne b$
gives, by comparing its two completions with tight diagonals,
\begin{align*}
 x_B(b)-x_B(a)&\ge AB[a,b]+BC[b,a]-AB[a,a]-BC[a,a],\\
 x_A(a)-x_A(b)&\ge AB[a,b]+AC[a,b]-AB[b,b]-AC[b,b].
\end{align*}
The cross terms must be retained for repeated-label $P$ teams. Any valid
core inequality $x_i\ge x_j+c$ also yields $H_i\ge H_j+c$ by using
a core point attaining $H_j$. A cycle of such difference inequalities
has total right-hand side at most zero; if the total is zero, every
inequality is tight at every core point. These are useful consequences,
but discarding the remaining team constraints need not describe the
entire core or its ideals. The compressed system is not claimed for
arbitrary dense supports.

\subsection{Linearly many incompatible ideal faces}
\begin{proposition}\label{anc:attainment}
In the family of Proposition~\ref{prop:family}, let $U$ be the players
outside its $\ell$ special teams. Any original core allocation attains
$H_i=3$ for at most one $i\in U$. Exactly $3(\ell+1)$ core allocations
are necessary and sufficient to attain all those unmatched ideals
collectively.
\end{proposition}
\begin{proof}
There are $\ell+1$ unmatched players in each sector. If player $(s,j)$
receives three at a core point, diagonal tightness makes the other two
coordinates at label $j$ zero. Complete their opposite-sector diagonal
pair of value one with any sector-$s$ player. Its team inequality forces
every sector-$s$ payoff to be at least one.
A second ideal three in another sector would likewise force every
payoff in that sector to be at least one, contradicting its zero payoff
at $j$. If two unmatched ideals in the same sector were attained, that
sector's $\ell+1$ unmatched players would receive at least
$(\ell+1)+4=\ell+5$ in total. But summing the $\ell$ disjoint
special-team constraints gives
\[
                  x(U)\le T-5\ell=\ell+3,
\]
a contradiction. Each unmatched ideal face is nonempty by the witnesses
in Proposition~\ref{prop:family}; choosing one witness per face attains
the lower bound $|U|=3(\ell+1)$ on the required number of points.
\end{proof}
A bounded collection of allocations that exactly attain the ideals
cannot therefore work uniformly over this family. This does not exclude
fractional protection obtained from nonmaximizing allocations. In
contrast, if one core point simultaneously attains every ideal of each
sector, the average of the three points already protects $H/3$ with
zero subsidy. Published side-optimality in the square supermodular class
with local monotonicity~\cite[Theorem~3.2]{Martinez2026} supplies such a
safe class under that theorem's additional hypotheses; neither the
side-optimality concept nor this elementary averaging observation is
claimed as a new contribution.

\subsection{A ring and the effect of weakening one pair}
For $b\ge2$, take $m=4b$ labels and initialize all three pair matrices
to the identity. Let the actual partial matching $P$ consist of
\[
 e_j=(j,(j+1)\bmod b,b+j),\qquad
 f_j=(2b+2j,2b+2j+1,j),\qquad 0\le j<b.
\]
For every displayed team set its off-diagonal pair values, in order
$AB,AC,BC$, to $(2,1,2)$; other off-diagonal values are zero.
Call the resulting market $R_b$. Call $R'_b$ the market obtained by
lowering the single entry $AB[0,1]$ from two to zero. The teams in $P$
are player-disjoint and have three distinct labels.

\begin{theorem}[An all-size perturbation family]\label{anc:ring}
Both $R_b$ and $R'_b$ have nonempty original cores and grand worth
$T=12b$. Their protected subsidies are
\[
                          s(R_b)=0,\qquad s(R'_b)=4b-2.
\]
Every ideal of $R'_b$ equals three. Each market's ideals are calculated
in its own original core and then frozen during protection.
\end{theorem}
\begin{proof}[The unmodified ring]
At each hub label $j<b$ pay $4/3,4/3,1/3$ in sectors $A,B,C$.
At each other label pay $7/3$ to its unique $P$-used player and $1/3$
to its other two players. Denote this vector by $y$. Every diagonal
team receives three, every $P$ team receives five, and $y(N)=12b$.
Two positive off-diagonal pairs in a team force it to be a $P$ team,
because $P$ is a partial matching. Any other nondiagonal team has at
most one such pair and at most one diagonal pair. For a pair of value
two, its endpoints receive at least $8/3$ and its third player at least
$1/3$, covering total worth at most three. For a pair of value one,
the endpoints receive at least $8/3$, covering total worth at most two.
With no off-diagonal pair, worth is at most one, again covered. Hence
$y$ covers every team and every coalition. Together with the diagonal
matching this proves grand worth $12b$ and core nonemptiness.

At any original core point, diagonal tightness and the completion
inequalities for $(j,j+1,j)$ and $(j,j+1,j+1)$, with hub indices taken
modulo $b$, give $x_B(j+1)\ge x_B(j)$ and $x_A(j)\ge x_A(j+1)$.
Around the ring these become equalities. Write the common hub payoffs
as $A,B,C$ with $A+B+C=3$, and put $c=\min_k x_C(k)$.
The ring-team inequality gives $x_C(b+j)\ge2+C$.
Completing the diagonal $AB$ pair at $b+j$ with a sector-$C$ minimum
gives
\[
       1-c\le x_A(b+j)+x_B(b+j)=3-x_C(b+j)\le1-C.
\]
Thus $c=C$, $x_C(b+j)=2+c$, and the other two payoffs there sum to
$1-c$. In particular $0\le c\le1$ and those two coordinates are at
most one.

For attached labels $u=2b+2j$, $v=2b+2j+1$, let
$a=\min_kx_A(k)$ and $d=\min_kx_B(k)$. With $p=x_A(u)$ and
$q=x_B(v)$, the attached-team inequality is $p+q+c\ge5$.
Diagonal tightness gives $q\le3-a-c$ and $p\le3-d-c$.
Hence $p\ge2+a$, $q\ge2+d$, and each of the other coordinates at
$u,v$ is at most one. All these estimates hold on the entire core.
Thus unmatched players and hub-$C$ players have ideal at most one;
all other players have ideal at most three by diagonal tightness.
Vector $y$ pays the former group $1/3$ and every other player at least
one, so it is already an efficient protected allocation.
\end{proof}

\begin{proof}[The broken ring]
Put $z=b$; the altered team is $e_0=(0,1,z)$, of worth three.
The other $2b-1$ teams in $P$, called hard teams, have worth five.
Removing $e_0$ from the label--team incidence graph leaves $z$ isolated
and a tree on the other $4b-1$ labels: the remaining ring is a path,
and each $f_j$ attaches two leaves to a hub. Call this the main tree.
Root it at any label with any sector color. Traversing outward, color
each newly visited label by its sector position in the traversed team.
Every hard team then has at least two endpoints with their designated
colors, so every pair in it has at least one correctly colored endpoint.
Let $c^A$ be the coloring rooted at $0$ with color $A$, and $c^B$
the coloring rooted at $1$ with color $B$. The ring path gives
\[
                c^A(0)=c^A(1)=A,\qquad c^B(0)=c^B(1)=B.
\]

To attain ideal three at coordinate $(k,v)$ with $v\ne z$, root the
main tree at $v$ with color $k$, extend its coloring by $c(z)=C$,
and pay
\[
                        x_{s,j}=\one_{\{s=k\}}+2\one_{\{c(j)=s\}}.
\]
For a target at $z$, pay three to its target sector there and zero to
the other two sectors. On every main-tree label use, respectively,
\begin{align*}
 \text{target }A:\quad x_{s,j}&=\one_{\{s=A\}}
                         +\one_{\{c^A(j)=s\}}+\one_{\{c^B(j)=s\}},\\
 \text{target }B:\quad x_{s,j}&=\one_{\{s=B\}}+2\one_{\{c^B(j)=s\}},\\
 \text{target }C:\quad x_{s,j}&=\one_{\{s=C\}}+2\one_{\{c^A(j)=s\}}.
\end{align*}
Every label has total three, the target receives three, and every player
in the target sector receives at least one.

We check all team inequalities. A hard team receives at least four
coloring units and one base unit, covering five. Endpoints of a hard
pair receive at least two coloring units, covering its value at most
two. If a team containing that pair also has a diagonal pair, its three
labels lie on the main tree, so it also receives a base unit and covers
worth at most three.

The only positive soft pairs are $AC(0,z)=1$ and $BC(1,z)=2$.
For a target on the main tree, $C_z$ has two coloring units and every
full team has a base unit, so any team containing a soft pair receives
at least three. For target $C$ at $z$, $C_z=3$ gives the same coverage.
For target $B$ at $z$, labels $0,1,z$ all have vector $(0,3,0)$.
The soft $BC$ pair receives three; the soft $AC$ pair receives zero
but every $B$ coordinate is at least one, and if there is also a diagonal
pair its $B$ label is $0$ or $z$, where it receives three. For target
$A$ at $z$, labels $0,1$ have vector $(2,1,0)$ and $z$ has vector
$(3,0,0)$. The soft $AC$ pair receives two. The soft $BC$ pair
receives one and every $A$ coordinate is at least one; with an additional
diagonal pair, the $A$ label is $1$ or $z$ and the team receives at
least three. The soft team itself receives three in each case.

Two positive off-diagonal pairs still force a $P$ team. Every other
nondiagonal team has at most one such pair and at most one diagonal
pair, so the preceding checks exhaust these cases. Without off-diagonal
pairs its worth is at most one, covered by the target-sector lower
bound. Diagonal teams receive exactly three. Thus every displayed
witness is a total-$12b$ core point. The diagonal matching has this
value, and its tightness bounds every ideal by three. Separate witnesses
attain three at every coordinate.

For the exact subsidy, start with one per player and add two to the
$A$ player of each hard team. These players are distinct. Hard teams
receive five, the soft team receives three, and all other teams have
worth at most three. This covers all coalitions and the floors, with
total $12b+2(2b-1)$. Conversely, the disjoint teams in $P$ have total
worth $5(2b-1)+3=10b-2$, and their $6b$ unmatched players require one
each. The resulting lower bound on protected total is $16b-2$, matching
the construction.
\end{proof}

\begin{corollary}[No modulus uniform over market size]\label{anc:modulus}
Scale $R_b,R'_b$ by $1/(12b)$, giving grand worth one in both games,
and put $\delta_b=1/(6b)$. Their pair-matrix entrywise $\ell_1$
distance and their characteristic-function supremum distance both equal
$\delta_b$, whereas their protected subsidies differ by
$1/3-\delta_b$. Consequently no function $\omega(\delta)\to0$ as
$\delta\to0$ uniformly bounds this subsidy difference over all finite
market sizes in either distance.
\end{corollary}
\begin{proof}
Only one pair changes, by $\delta_b$. An integral matching uses that
pair at most once, so every coalition worth changes by an amount in
$[0,\delta_b]$. The three-player coalition $e_0$ changes by exactly
$\delta_b$. Positive scaling scales the core, ideals and subsidy,
so Theorem~\ref{anc:ring} gives the asserted subsidy difference.
\end{proof}
This is a failure of a size-independent modulus, not discontinuity at a
fixed market size. It also shows why removing a small pair value cannot
be charged a universal proportional cost when the original ideals change
with the market. Instability of selected assignment-core allocations
has established precedents~\cite[Example~3]{anc:KumabeYoshida2024}; the
scalar protected-subsidy statement and its priority require their own
comparison.

\subsection{An optimal obstruction that must contain cycles}
\begin{proposition}[A finite certificate]\label{anc:cycles}
There is a three-sector market with $m=11$ labels, grand worth $T=1$,
nonempty entire core and subsidy $76/303$, whose unique optimal grand
matching is diagonal. Every optimal team-form subsidy dual contains six
specified disjoint teams, each with multiplier one. Their connected
label--team incidence graph has cycle rank two. In particular, no
forest-supported optimal team-form dual exists, including fractional
duals or a change to another optimal grand matching.
\end{proposition}
The statement is computer-assisted through a supplied rational
certificate, not an asymptotic conclusion from experiments. We specify
the market and the finite proof obligations explicitly. All entries in
the following tables are divided by $5252$; every omitted pair entry
is zero. The diagonal $AB,AC,BC$ entries are
\[
\begin{array}{c|rrr@{\qquad}c|rrr}
 j&AB&AC&BC&j&AB&AC&BC\\\hline
 0&108&234&144&6&153&153&144\\
 1&252&144&90 &7&189&144&153\\
 2&144&198&198&8&144&144&198\\
 3&90&108&108 &9&198&144&198\\
 4&154&144&148&10&198&144&198\\
 5&108&234&144&&&&
\end{array}
\]
and the six off-diagonal teams $P$ have pair entries
\[
\begin{array}{c|rrr}
 (A,B,C)\text{ labels}&AB&AC&BC\\\hline
 (1,0,2)&252&288&252\\
 (3,4,1)&158&144&234\\
 (0,5,3)&252&198&108\\
 (6,3,7)&153&297&153\\
 (8,9,5)&288&234&306\\
 (7,10,8)&333&198&342
\end{array}
\]

\begin{proof}[Certificate proof]
The record \texttt{CYCLIC\_FRONTIER\_6\_2\_0} in the supplied file
\texttt{CYCLIC\_FRONTIER\_CERTIFICATES.json} contains all $33$ exact
ideals, a separate original-core point attaining each, their matching
upper certificates, and a protected allocation $y$ with a matching dual
of total $1+76/303$. Each core witness has nonnegative coordinates,
total one and all full-team inequalities. For the upper certificate of coordinate $i$, the stored multipliers
$\lambda_t\ge0$ and the efficiency multiplier $\nu$ satisfy
\[
 \nu-\sum_{t\ni j}\lambda_t\ge\one_{\{j=i\}}
 \quad(j\in N),\qquad
 \nu-\sum_t\lambda_tv(t)=H_i.
\]
Since $x\ge0$, $x(N)=1$ and $x(t)\ge v(t)$, these equations imply
$x_i\le\nu-\sum_t\lambda_tv(t)=H_i$. The multiplier on efficiency
is unrestricted; no unsupported joint ideal attainment is assumed. The protected vector satisfies all
floors and full-team inequalities, while the dual multipliers satisfy
the player-load bounds and the displayed objective. These exact finite
checks prove $T=1$, all declared ideals and the subsidy value; they do
not substitute a numerical solver tolerance for an inequality.

Average the $33$ original-core witnesses. Direct rational checks show
that the resulting point $x$ is strictly positive and its only tight
full teams are the eleven diagonals and $(1,0,2),(7,10,8)$.
A grand matching of worth one must cover every player and use only
$x$-tight teams: its value deficit is the sum of team slacks and
unmatched positive payoffs. At each label there is a player appearing
in no other tight team. Sectors
\[
                       A,B,A,A,A,A,A,B,A,A,A
\]
provide such a player at labels $0,\ldots,10$, respectively.
Their diagonal teams are all forced, proving uniqueness. This reasoning
also rules out another fractional perfect optimum.

For the protected problem use full-team multipliers $\lambda_t\ge0$
and floor multipliers $\mu_i\ge0$, with
$\sum_{t\ni i}\lambda_t+\mu_i\le1$. The primal--dual difference
at the certified $y$ is
\begin{align*}
 &\sum_t\lambda_t\bigl(y(t)-v(t)\bigr)
 +\sum_i\mu_i\bigl(y_i-H_i/3\bigr)\\
 &\hspace{25mm}+\sum_i y_i
                  \left(1-\sum_{t\ni i}\lambda_t-\mu_i\right).
\end{align*}
All terms are nonnegative. The certificate has $y_i>0$ for every
player, and exactly twelve full-team constraints are tight. At any
optimal dual, non-tight teams have multiplier zero, strict-floor
players have $\mu_i=0$, and every player load is one. In order, the
following strict-floor players each occur in exactly one still
available tight team:
\[
\begin{array}{c|c}
 \text{player}&\text{forced team}\\ \hline
 A_0&(0,5,3)\\
 A_1&(1,0,2)\\
 B_9&(8,9,5)\\
 B_{10}&(7,10,8)\\
 C_1&(3,4,1)\\
 C_7&(6,3,7)
\end{array}
\]
Its multiplier is therefore one. All teams sharing any of its players
then have multiplier zero, justifying the exclusions before the next
step. The exact checker verifies each strict inequality and each unique
choice, rather than trusting the listed dual. Thus every optimal dual
contains all six teams with multiplier one. Their incidence graph is
connected with $18$ edges and $17$ vertices, so its cycle rank is
$18-17+1=2$. Adding further positive teams cannot remove these cycles.
\end{proof}
The certificate and checker are retained in
\nolinkurl{structural_classification_2026-09-06/sharp_subsidy/}.
The solver-free command
\texttt{python3 -B -O verify\_cyclic\_necessity.py}
reconstructs the market, verifies the original ideals and primal--dual
value, checks all strict/tight sets and forcing steps, and rejects a
false forcing step. Its direct dependency is
\texttt{verify\_cyclic\_frontier.py}, with its retained checker imports
and the named JSON certificate; run from the complete supplied directory.
This result rules out a forest-only optimal-dual restriction on the
same market. It does not exclude a transformation that changes the
market while correctly tracking all changes to its original ideals.

\paragraph{Reproducibility and claim boundaries.}
The same retained directory contains separate exact checkers for
Proposition~\ref{anc:compression}, Proposition~\ref{anc:attainment}
and Theorem~\ref{anc:ring}: respectively
\texttt{verify\_sector\_min\_core.py},
\texttt{verify\_coattainability.py}, and
\texttt{verify\_broken\_ring.py} together with \texttt{ring\_zero.py}.
They have been run with Python assertions disabled. The all-size
statements rely on the written proofs; the cyclic-optimum proposition
is expressly a finite certificate proof. Historical priority and
independent human specialist assessment remain open.
